\documentclass[11pt]{article}
\usepackage[T1]{fontenc}
\usepackage{lmodern}
\usepackage[margin=1in]{geometry}
\usepackage{amsmath,amssymb,amsthm,booktabs,array,longtable}
\usepackage[hidelinks]{hyperref}
\usepackage{url}
\Urlmuskip=0mu plus 2mu
\setlength{\emergencystretch}{3em}
\newtheorem{proposition}{Proposition}[section]
\newtheorem{definition}[proposition]{Definition}
\newtheorem{remark}[proposition]{Remark}
\newcommand{\F}{\mathbb F}
\newcommand{\Z}{\mathbb Z}
\newcommand{\art}[1]{\nolinkurl{#1}}
\newcommand{\status}[1]{\textbf{Status: #1.}}
\title{Scoped Boundaries for Elliptic-Curve Index Calculus:\\Counterexamples, Conservation, and Oracle Models}
\author{Anonymous}
\date{Research draft, 18 October 2026}
\begin{document}
\maketitle
\begin{abstract}
We assemble six proposed negative results from an existing elliptic-curve
index-calculus research corpus, preserving their mathematical objects and
quantifiers. A universal comparison between true and Boolean-reduced first
fall degrees has explicit small counterexamples, but its interpretation depends
on the degree convention and the placement of field equations. A measured
failure of degree-four mutant closure does not by itself refute Semaev's F4
assumption. Table-based decomposition admits a generic simulation under
explicit restrictions; connecting a concrete implementation to that model is
an additional obligation. Mean decomposition multiplicity obeys an exact
finite-fibre counting identity, which already has a Lean formalization in the
corpus. The signed-sum description of the ordered Semaev cover gives elementary
abelian monodromy and a restricted specialization law. Isogeny-map oracles are
linear in the generic model; coordinate-aware oracle elimination instead
requires transported challenges and an algebraic-query restriction. This is a
source-bound synthesis and proof inventory. It introduces no scientific trials,
claims no new theorem or novelty, and makes no deployed-curve security claim.
\end{abstract}
\tableofcontents
\section{Scope, sources, and the meaning of a negative result}

An index-calculus proposal can fail for different reasons. Its counting
observable can be fixed identically; its proposed generic oracle can already
be simulated; an inference between degree notions can have a counterexample;
or a finite solver experiment can miss its registered threshold. These are
different conclusions. In particular, none establishes the nonexistence of a
future coordinate-based algorithm for the elliptic-curve discrete logarithm.

The six subjects here are: Nagao's first-fall comparison; the interpretation
of degree-four experiments relative to Semaev's assumption; generic simulation
of compact-orbit tables; conservation of mean decomposition multiplicity;
monodromy of the ordered summation cover; and isogeny-map oracles. Statements
below are restatements or elementary explanations of identified existing
sources. A proposition label does not assert originality or substitute for
independent review. The source map and outstanding obligations are part of
the draft, rather than qualifications to be removed before publication.

The corpus already contains an ECDLP programme manuscript in
\art{crypto:papers/ecdlp-index-calculus-program/} and a conservation and
specialization manuscript in
\art{crypto-autoresearcher:papers/semaev-conservation-specialization/}.
Both identify themselves as unsubmitted drafts. This companion is therefore
not the first ECDLP manuscript in the programme. Its purpose is to reconcile
precise scope, including contradictions between historical summaries and
later source tests. It is not a publication or submission receipt.

\paragraph{Evidence vocabulary.}
A finite algebraic counterexample can disprove a universal statement, even
when its variables are few. A converged closure calculation describes its
specific representation and degree convention. An asymptotic generic lower
bound describes an oracle model, not all computations on coordinates.
An exact double count describes multiplicity, not search cost. An
extrapolation of a cost model is not a measurement. We keep these distinctions
throughout.

\begin{center}
\fbox{\begin{minipage}{0.88\linewidth}
\textbf{Claim chain to preserve}\par
Mathematical object and quantifiers
$\longrightarrow$ statement or frozen protocol
$\longrightarrow$ proof or retained computation
$\longrightarrow$ independent verification
$\longrightarrow$ the conclusion with its original scope.\par
A missing arrow is an open obligation; a source title cannot replace it.
\end{minipage}}
\end{center}

\section{First fall comparisons and degree-four measurements}
\subsection{A comparison of two degree conventions}

Nagao's Definitions 5 and 6 distinguish a degree in
$\F_p[X_1,\ldots,X_N]$ from a degree after reduction modulo the field
equations. Lemma 4 compares $d_F(F)$ to $d'_F(F\cup S_{\rm fe})$, with
$S_{\rm fe}=\{X_i^p-X_i\}$ \cite{nagao2015}. The frozen extraction explicitly
places the field equations outside the true-side generator set in that
lemma. The programme's current certificate code makes the intended
comparison convention explicit \cite{nagaoCertificate}.

For the equality convention, a fall at $D$ has products $g_i f_i$ whose
maximum degree is exactly $D$, and whose nonzero sum has degree strictly
less than $D$. In the fake quantity every product and the sum is first
reduced modulo $S_{\rm fe}$. The requirement that the maximum is attained
must be checked; a low-degree element in a span alone is insufficient.

\begin{proposition}[Existing counterexample, equality convention]
Over $\F_2$, the universal inequality $d_F\leq d'_F$ fails. Two recorded
witnesses are
\[
 F_A=\{X^2Y,\;XY+X\},\qquad
 F_B=\{X_1X_2+X_3,\;X_1X_3+X_2\}.
\]
The current certificate tests assign fake degree $2$ and true degree $3$
to each witness, for either true-side placement of the field equations.
\end{proposition}

The witnesses and exact test assertions are in
\art{crypto-autoresearcher:tools/lemma4_inside_form.py} and its test module.
For $F_B$, take $g_1=1$ and $g_2=X_1+1$. Boolean reduction gives
\[
 \overline{f_1+(X_1+1)f_2}=X_2+X_3,
 \quad \deg\overline{f_1}=\deg\overline{(X_1+1)f_2}=2.
\]
The ordinary products attain degree $3$, not $2$. The tests separately
check all degree-two scalar combinations when field equations are
adjoined, and exhibit a degree-three fall. These checks are finite
certificates for the stated polynomial systems, not elliptic-curve
instances or measurements of an F4 run.

\paragraph{The printed inequality reading is a separate axis.}
The source prints a $\geq$ form of the degree condition. Under the
programme's padding analysis that reading collapses the quantity to one
plus the least degree of a nonzero ideal element. The current test class
\art{TestTheGeReadingVoidsOnlyOnePlacement} distinguishes the true-side
placement as follows, for both witnesses:
\begin{center}
\begin{tabular}{llll}
\toprule
Condition reading & True-side generators & True & Fake \\
\midrule
Equality & $F$ & 3 & 2 \\
Equality & $F\cup S_{\rm fe}$ & 3 & 2 \\
Literal $\geq$ & $F$ & 3 & 2 \\
Literal $\geq$ & $F\cup S_{\rm fe}$ & 2 & 2 \\
\bottomrule
\end{tabular}
\end{center}
Thus the assertion that the literal reading always voids the
counterexample is too broad. The historical finding
\art{KN-FIND-936151} contains that broader wording; the later tests
explicitly correct it. Both are retained in the source inventory. The
frozen extracted Lemma 4 names $F$ on the true side, but a final
publication still owes a fresh PDF-level transcription of both definitions
and their intended reading. These sources do not prove or disprove
Nagao's Proposition 5 by another argument, nor a complexity conclusion
for ECDLP. The programme itself withholds the broader derived claim.

\subsection{What the binary degree measurements actually establish}

Semaev's Assumption 1 concerns $d_{F4}$ for the Boolean system equivalent
to his chained equations, with subspace dimension $k=\lceil n/m\rceil$
\cite{semaev2015}. It is distinct from Nagao's separately numbered
Assumption 1, from first fall degree, from a one-pass Macaulay rank, and
from a bounded-degree mutant closure. Confusing these quantities can
turn a valid computation into an invalid refutation.

The closure-at-scale study \cite{closureScale} reports eight unsatisfiable
$m=3$ draws at $(n,\ell)=(19,7)$ and $(21,7)$ for which the complete
degree-four mutant closure saturates without containing $1$. It also
reports degree-four refutation on its other listed cells, including the
$m=4$ cells through $\ell=4$. This is a failure of a uniform
constant-four closure model in that representation. It does not follow
merely from a timeout: the report says these particular closures
converged. A new verifier would still need the retained instance, closure
basis, closure-stability checks and non-membership evidence before
certifying an individual negative instance afresh.

\status{The broad claim that these records refute Semaev's Assumption 1
is not established.} The programme's own correction
\art{KN-FIND-9cba50} records that an earlier ladder used
Frobenius-invariant subspaces of dimension $\operatorname{ord}_n(2)$,
which need not equal $\lceil n/m\rceil$, and that its degree meters were
not interchangeable with $d_{F4}$. Conversely, the two $m=3$ cells just
listed do have $\ell=\lceil n/m\rceil$; calling them off-diagonal does
not fix the metric issue. A verified bridge from saturated mutant
closure to the precise F4 quantity remains necessary in either case.

There is also positive counterevidence against a naive degree-meter
identification. At the separately archived $n=17,m=2,\ell=9$ cell,
plain degree-four Macaulay left 62 unsatisfiable targets unresolved;
the subsequent mutant closure refuted all 62 with checked witnesses
\cite{closureCorrection}. Therefore non-refutation by plain Macaulay
was not itself evidence that mutant refutation degree exceeded four.
No statement about cryptographic-size degrees follows from this finite
example. Likewise, fitting either of two surviving degree-growth models
does not logically refute every possible $m=4$ degree surface outside
the measured cells.

\section{Compact-orbit tables: a conditional generic closure}

The programme's step-back note classifies compact-orbit decomposition
through enumeration, pair tables and group-law root extraction as generic
\cite{compactSource}. The useful mathematical content is a simulation
statement with explicit admissible operations.

\begin{proposition}[Simulation schema, conditional on the interface]
Suppose a prime-order-group computation constructs its advice through
generic group operations and public known scalar actions. Its online
phase uses only those operations, equality, and lookup against advice
sets; representation-dependent indexing is required to be bookkeeping
with no additional mathematical predicate. Then the table-decomposition
computation can be simulated in the same generic interface, charging its
advice and all simulated operations.
\end{proposition}

The simulation tracks each available point as $[a]G+[b]Q$. Addition,
negation and multiplication by a public scalar update $(a,b)$; a table
entry is another stored point, and membership against the table is an
equality query or a generic lookup whose implementation cost is charged.
For a fixed known factor-base point $F$, extraction through
$S_3(x_Q,x_F,X)$ can instead compute the group elements $Q+F,Q-F$.
This identity does not supply a generic method for deciding arbitrary
low-degree coordinate restrictions on unknown summands. The latter
operation is precisely what an algebraic decomposition method must
exploit, and is outside the proposition's interface.

The generic discrete-log lower bound \cite{shoup1997} then applies to a
computation that meets these hypotheses. With target-independent advice,
Corrigan-Gibbs and Kogan give an appropriate preprocessing tradeoff
\cite{cgk2018}; its frequently quoted schematic form is
$S T^2=\widetilde\Omega(\varepsilon r)$, with success probability
$\varepsilon$ and the original paper's units and parameter regime.
An automorphism quotient needs its own model and orbit accounting before
replacing $r$ by an effective orbit count. Neither schematic inequality
proves that one concrete implementation beats every other implementation
in wall time or at equal bytes. Such a constant comparison needs paired
measurements and an equal accounting interval.

\paragraph{The concrete bridge is an obligation.}
A factor base selected by coordinates is not automatically generated
by a generic interface. A lookup key can be mere indexing, whereas a
field trace, a subspace test or a coordinate-dependent branch can supply
additional information. One must inspect the frozen construction and
all online predicates, state whether and how their outputs are simulated,
and account for any target-independent coordinate-selected advice. This
draft does not certify every current compact-pipeline path against a
new simulator. In particular, the source's historical statement that
all table methods are dominated by preprocessing rho at every constant
must not be substituted for a same-input comparison.

The existing Lean file \art{GenericInvariance.lean} proves transcript
invariance for a more restrictive interface: \art{Alg.next} sees an
operation count, prior equality answers and a random tape, but no point
representation. Its source comment explicitly warns that it does not
cover the representation-hashing rho and kangaroo implementations that
ecbench measures. It proves neither the current compact pipeline's
simulation bridge nor the probabilistic preprocessing lower bound.
The two model boundaries must remain visible when using that file.

\section{Conservation is an exact statement about multiplicity}

The conservation identity is the least assumption-sensitive result in
this inventory \cite{yieldSource}. It is a finite-fibre count.

\begin{proposition}[Finite-pattern conservation]
Let $M$ and $G$ be finite sets and $s:M\to G$ a function. Put
$c(g)=|s^{-1}(g)|$. Then
\[
 \sum_{g\in G}c(g)=|M|.
\]
If $G$ is nonempty and the target is uniform, its mean multiplicity is
$|M|/|G|$.
\end{proposition}
\begin{proof}
Each element of $M$ belongs to exactly one fibre of $s$. The fibres
partition $M$, so their cardinalities sum to $|M|$. Division by the
number of targets gives the stated uniform mean.
\end{proof}

For a finite abelian group $G$, a set $D$ of $B$ distinct elements and
$m\geq1$, choose $M$ according to the declared counting convention:
\begin{align*}
 \text{ordered tuples with repetition:}\quad &|M|=B^m,\\
 \text{unordered multisets with repetition:}\quad
 &|M|=\binom{B+m-1}{m},\\
 \text{one element from each typed set }D_i:\quad
 &|M|=\prod_{i=1}^m |D_i|.
\end{align*}
The sum is well defined for multisets because the group is abelian.
Comparing two conventions without changing the denominator is not a
violation of conservation. Signed decompositions similarly require
specifying whether signs are separate labels or whether $D\cup(-D)$ is
a deduplicated set.

Coverage $p=|\{g:c(g)>0\}|/|G|$ is not the mean $\mu$.
Exactly,
\[
 p=\mu-\frac1{|G|}\sum_g(c(g)-1)_+\leq\min(1,\mu).
\]
When $\mu\leq1$, equality $p=\mu$ requires every count to be zero or one;
when $\mu>1$, equality with the bound $1$ requires coverage of every
target and permits repeated decompositions. Conservation therefore
leaves coverage, concentration and useful relation rank free. A claimed
universal numerical bound on a structured base's gain over a random base
would additionally require a model or bound for that random comparator;
it does not follow from this double count alone.

\paragraph{What is not conserved.}
The identity is not about the cost of finding a decomposition, the rank
of resulting rows, a nonuniform target law, or the mean over a selected
proper subset of $G$. A learned filter can redistribute work across
targets without contradicting it. A target-dependent base is another
experiment and must charge its selection. These are genuine open axes,
not exceptions to the stated identity.

The corpus already formalizes the general finite-pattern core in
\art{FactorBaseDesigns.sum_count}. It uses a list domain and a
nonduplicated target list covering its image, so domain repetitions
count as labelled multiplicities. The ordinary tuple or multiset
cardinality formula is a separate instantiation. Section~\ref{sec:formal}
records that boundary precisely.

\section{The ordered summation cover and its specialization}

The relevant cover must be named before giving its monodromy. Let $E/K$
be an elliptic curve, let $\pi:E\to\mathbb P^1$ be the quotient by
inversion, put $n=m-1$ with $m\geq3$, and take independent ordered
transcendentals $t_1,\ldots,t_n$. The corpus's theorem concerns
\[
 S_m(t_1,\ldots,t_n,T)\quad\text{over }L=K(t_1,\ldots,t_n),
\]
not an unspecified symmetric parameter quotient or every specialization
\cite{monodromySource}.

\begin{proposition}[Existing signed-sum structure statement]
Assume the quotient map $\pi$ is separable of degree two with deck
involution $[-1]\ne1$, and the signed sums are distinct at the generic
point up to simultaneous inversion. The root field and splitting field
are the fixed field of diagonal inversion in $K(E^n)$, and their group
over $L$ is
\[
 (\Z/2\Z)^n/\langle(1,\ldots,1)\rangle
       \simeq(\Z/2\Z)^{m-2},
\]
acting regularly on the $2^{m-2}$ roots.
\end{proposition}

The existing derivation obtains the product deck group from the
independent quotient maps. Its roots are
$x(\epsilon_1P_1+\cdots+\epsilon_nP_n)$. Diagonal inversion fixes all
roots. A further generic coincidence would force a nonempty proper
partial signed sum into $E[2]$; the preimage of that finite group scheme
under the corresponding surjective sum map is proper. This gives the
stabilizer exactly, and abelianness makes the single-root field the
splitting field. The corpus's longer manuscript also argues that the
constant field is unchanged, hence arithmetic and geometric monodromy
agree. A specialist review of the quotient and characteristic-two
scheme argument remains on its publication checklist; this synthesis
adds no independent proof review.

Over a finite field, exclude branch points, infinite signed sums, degree
drops and coincident roots. Frobenius on each degree-two fibre is a sign
bit. In the regular quotient action its image has order at most two.
Consequently a good specialization is either completely split or a
product of equal-degree irreducible quadratics. If all inputs are
abscissae of rational factor-base points, their sign bits are zero and
it is completely split. This last assertion also follows directly:
compute the signed sums of those rational points.

\paragraph{The exclusion this supplies.}
A factorization-type statistic is constant on this good rational
factor-base locus. Such a statistic cannot distinguish relation-rate
anomalies there. It does not follow that every statistic or algebraic
representation is constant, or that finding the other summands for a
fixed target is cheap. Complete splitting solves for the last coordinate
when the others are given; relation search asks a different question.

The older finding \art{KN-FIND-a1f3c2} had declared a broader open problem
closed. The later \art{KN-FIND-6edecd} explicitly declines that closure:
its separate toy corroboration addresses the full-symmetric-monodromy
premise, while the relation-rate and imprimitivity questions are not
closed by that empirical study. We therefore report the ordered-cover
structure and its stated exclusions, not the older universal closure
language. Permuting or symmetrizing the input coordinates changes the
base and can enlarge the monodromy; the elementary abelian description
must not be transferred to that different cover without a proof.

\section{Isogeny oracles: generic and coordinate-aware models}

\subsection{The generic statement}
The source's Theorem 15.1 works with independent random encodings of
order-$r$ subgroups on curves in one rational isogeny class, where $r$
is prime and coprime to the cofactor \cite{isogenySource}. The action of
each isogeny on chosen subgroup generators is supplied as a public
scalar. A homomorphism between these cyclic subgroups is either zero
or an isomorphism. Its answer on a point with formal logarithm $a+bx$
has formal logarithm $\lambda(a+bx)$ in the destination.

Thus addition, negation and map queries preserve linear formal
expressions. Unless two distinct expressions collide at the uniform
secret $x\in\F_r$, the encoded transcript supplies no equation for it.
For $L$ initial handles and $Q$ queries, a collision union bound has the
usual order $O((L+Q)^2/r)$, plus the probability of an uninformed final
guess. The source writes a particular $(Q+2)^2/r+1/r$ expression;
its exact constant requires fixing initial-handle and domain accounting.
We use only the model-dependent birthday-order implication here.

This is a statement about independently encoded subgroup operations,
not rational coordinate access. Known map scalars are part of its
hypotheses; the cost of constructing those scalars or a useful geometric
map cannot be inferred from isogeny degree alone. Pairing oracles and
$p$-adic lifts are outside this stated interface.

\subsection{Why coordinate-aware elimination needs more hypotheses}
The source's later algebraic-group-model analysis
\cite{agmIsogeny} permits coordinate computations but requires every
point submitted to an isogeny oracle to carry an algebraic
representation over received points. It eliminates the oracles given
a \emph{class-joint instance}: $(P_j,xP_j)$ with the same unknown $x$
on every touched curve. It also gives the converse reduction. Eliminating
those transported challenges further needs an explicit computable
transport, for example a known spanning tree of isogenies, with its
cost charged.

Accordingly, the defensible conclusion is relative equivalence to the
transported challenges under the declared query restriction. Without
that restriction, arbitrary coordinate-generated factor-base points can
be submitted: the oracle is then a transport service for new points,
and direct map evaluation has a degree-dependent cost. Without
computable transport, the class-joint inputs are additional information
rather than a free reduction to a single input.

The source also rejects ``isogenies change constants, never exponents''
as a general assertion: composite binary extension fields can contain
GHS-weak and non-weak members in one isogeny class. Its prime-field
exponent assertion is labelled a conjecture. Neither the generic theorem
nor the coordinate-aware relative reduction resolves it.

\subsection{Transported decomposition counts require the full pullback}
An elementary counting result usefully separates a transfer from a new
base choice. Let $\varphi:H\to H'$ be a homomorphism of finite abelian
groups, let $K=\ker\varphi$, choose $D'\subseteq\operatorname{im}\varphi$
and let $D=\varphi^{-1}(D')$. For ordered decompositions and $m\geq1$,
\[
 d_D(R)=|K|^{m-1}d_{D'}(\varphi R).
\]
For each codomain tuple, pick arbitrary preimages of its first $m-1$
entries. The last entry is uniquely determined by requiring their sum
to be $R$ and has the required image. There are exactly $|K|^{m-1}$
choices. This is the corrected finite-group form already used in the
programme manuscript \cite{transferSource}.

The image condition and full-pullback hypothesis matter. A codomain
point outside the image has no rational preimage; an arbitrarily chosen
subset of a fibre need not have $|K|$ lifts. The raw transfer note's
product-of-fibre-sizes display must be read tuplewise on the image, or
replaced by the formula above. Counts do not conserve arithmetic cost,
base membership degree, field of definition, or available endomorphisms.
A statement that every representation search can only lose does not
follow from this counting identity alone.

\section{The two requested Lean targets}\label{sec:formal}

\subsection{Yield conservation: a core theorem already exists}
The file
\art{crypto:research/factor_base_designs_20261015/lean/FactorBaseDesigns/Yield.lean}
defines \art{count D phi g} as the length of the filtered list of domain
entries mapping to $g$. Its theorem \art{sum_count} states, exactly in
mathematical notation,
\[
 \text{$Gs$ has no duplicates and contains every $\varphi(x)$ for $x\in D$}
 \quad\Longrightarrow\quad
 \sum_{g\in Gs}\operatorname{count}(D,\varphi,g)=\operatorname{length}(D).
\]
It requires decidable equality of the target type, not a group law.
\art{patterns_neutral} compares two equally long domains;
\art{restrict_le} proves that filtering the domain cannot increase
any target's multiplicity. The committed axiom log lists
\art{propext} and \art{Quot.sound} for these three declarations. A
committed build log reports a successful build. The present campaign also
rebuilt this nine-job project with its pinned Lean 4.34.1 toolchain and
reran its axiom audit; fresh logs are retained under
\art{crypto:research/ideas_execution_20261018/runs/}.

The campaign's new \art{IdeasExecution.OrderedYield} module enumerates the
ordered Cartesian power and proves \art{tuples_length},
\art{ordered_total_yield} and \art{ordered_patterns_neutral}. Its build
passes; its separate axiom audit lists only \art{propext} and
\art{Quot.sound}. This is an application bridge for ordered list-domain
multiplicity, not a new conservation principle.
For the multiset convention one still needs the enumeration length
$\binom{B+m-1}{m}$ without duplicate multiset entries. Prove that the
sum map and full target enumeration meet \art{sum_count}'s hypotheses.
For a mean in $\mathbb Q$, require a nonempty target set before dividing.
Those remaining bridges instantiate the existing theorem rather than create
a new conservation principle.

\subsection{Compact-pipeline genericity: an interface bridge is missing}
The exact theorem \art{GenericInvariance.run_map} in
\art{crypto:research/conductor_gap_cm_pairs_20261010/lean/GenericInvariance.lean}
states that an injective map preserving the declared multiplication and
inversion operations commutes with every bounded run of its handle-only
algorithm. \art{dlog_invariant} specializes it to a two-handle DLP input.
The source contains no elliptic-curve construction and no probabilistic
DLP lower bound. Its explicit representation-hashing limitation is
binding, as discussed above.
The campaign also rebuilt this four-job project and reran its axiom audit;
successful checking does not extend its interface to representation hashing.

A faithful target for the compact pipeline has three layers:
\begin{enumerate}
\item Specify a concrete fragment with generically constructed advice,
known scalar actions, membership through stored equality, and all index
operations exposed. Prove an abstraction relation between concrete
states and symbolic states across every executed instruction.
\item Prove a transcript simulation, including target-independent setup,
lookup/multiplicity semantics, exceptional cases and the exact resource
counters. Distinguish random-label hashing from the stricter
handle-only theorem rather than extending that theorem by prose.
\item Import or separately formalize the appropriate generic
preprocessing lower bound with its success probability, advice units,
automorphism assumptions and parameter range. It is not a corollary of
injective transport invariance alone.
\end{enumerate}
No inspected Lean file closes all three layers for the current compact
pipeline. A machine check of a small rewritten toy algorithm would not
close the concrete-code correspondence. This target therefore remains
open. The draft does not insert an unproved axiom or a \art{sorry} and
call that a formalization.

\section{Proof map, reproduction, and release obligations}

The machine-readable \art{source-map.json} binds the exact repository
revisions and Git objects of the principal sources. \art{PROOF-MAP.md}
records theorem-to-source and formalization obligations in a form that
can be reviewed without reading a build log. No historical result is
re-scored or promoted here.

The following commands are reproduction pointers, not receipts for
execution by this manuscript. Run them from the stated repository and
retain new logs without overwriting prior artifacts:
\begin{itemize}
\item In crypto-autoresearcher,
\art{python3 -m unittest discover -s tools -p test_lemma4_inside_form.py}
checks the current certificate convention and placement table. Read the
frozen primary PDF before drawing a conclusion about its intended text.
\item In crypto, the closure-at-scale report binds its preregistration,
raw rows and binary digests. A replay must reconstruct the same instance,
check unsatisfiability independently, and verify saturated closure
non-membership; a fresh generic solver timeout is not such a replay.
\item In crypto's factor-base Lean project,
\art{lake build}, \art{lake env lean AxiomAudit.lean}, and
\art{lake env leanchecker --fresh FactorBaseDesigns}
are the source's reproduction commands. Preserve both the toolchain
identity and fresh outputs.
\item The existing conservation/specialization paper identifies its
\art{EXP-SMON-e5cbe6} verification code and its deterministic retained
run. Reproduce into a new output location before reporting new numerical
verification. Finite samples remain corroboration of the derivation,
not a substitute for its proof.
\end{itemize}

\paragraph{Before any submission or theorem promotion.}
Resolve the primary-text transcription and current-test disagreement
for Nagao; bind a valid degree notion and instance family for any
Semaev refutation; audit the compact implementation's exact generic
interface; instantiate and freshly build the yield theorem; obtain a
specialist check of the ordered-cover and characteristic-two arguments;
and fix the exact isogeny-oracle domain/initial-handle accounting.
Independently inspect every source and bibliography entry at the
submission revision. Novelty, authorship, external publication and
independent review are not established by this draft. Publication is a
separate action with its own receipt.

The source inspection was authorized as item 19 of a twenty-item
execution request. The Conductor scope check was attempted before
edits and returned \art{conductor: command not found}. The session had
no knowledge-retrieval tool and used the indexed primary repository
artifacts instead. These are disclosed tooling limits, not findings
about elliptic curves. Drafting this manuscript partially advances
item 19; it does not complete its publication and formalization gates.

\bibliographystyle{plain}
\bibliography{refs}
\end{document}
